% Source: Independent numeric contract; CP-Algorithms basic-geometry; KACTL Point/Angle conventions.
\subsection{精度约定与退化情况}
\label{geo:precision}
\subsubsection*{原理}

先判断题目需要的是\textbf{组合关系}还是\textbf{实数构造}。整数输入下，共线、左右侧、线段相交、凸包、点在多边形内等优先用整数精确判断；只有需要长度、角度、交点、圆心时才转为浮点。整数乘法必须在乘之前提升到 128 位，不能先用 long long 溢出再转换。

本章统一使用以下约定：
\begin{itemize}
 \item 整数点为 IPoint，通常要求输入坐标绝对值不超过 $10^{12}$、点数不超过 $2\times10^5$。点坐标与差向量用 long long，叉积、距离平方和面积累加用 128 位；点坐标的加减也必须不溢出。后续重复求和或放大坐标时要重新估计范围。
 \item 实数为 long double，记作 ld；GEO\_EPS=$10^{-10}$ 是\textbf{长度的绝对容差}，GEO\_ANG=$10^{-18}$ 是单位方向之间的平行判定阈值。直线方向先归一化，side 的单位因此是长度；面积、距离平方不能不加换算就与长度容差比较。
 \item 浮点输入和计算中间量均须有限；通常将有效长度归一化到约 $[10^{-6},10^6]$，避免同时混用极大坐标和极微小间距。若两种拓扑关系只差不到容差，应按题目误差要求处理，浮点分支不承担精确判定。long double 的实际有效位数与编译器有关，可查 \InlineCode{numeric_limits<ld>::digits}；本地与 Linux 的 GCC 通常比 double 更宽，MSVC 则不一定。
 \item sort、set 的比较必须满足严格弱序。整数排序使用精确字典序／极角序；实数点的字典序也直接比较坐标，\textbf{不把带 EPS 的相等或大小关系放进排序器}。
 \item 多边形和点集按 0 下标，循环边为 i 到 (i+1) mod n，末尾不再重复首点；逆时针有向面积为正。有向直线的左侧为保留侧，圆半径非负，角度统一用弧度。
\end{itemize}
长度尽量用 hypot，比较距离时尽量比较平方；反三角函数输入夹到 [-1,1]。sqrt 前的非负截断只能消除已确认合法的公式的舍入误差，不能用它掩盖真正无交点的情况。接近平行的交点、接近共线的外心本身就敏感，应先平移、合理缩放，必要时用精确谓词或更高精度重新求解。

默认浮点阈值与校验以 \InlineCode{numeric_limits<ld>::digits>=64} 为基准；若平台的 long double 与 double 相同，须重新评估精度，不能只换类型名后照搬阈值。使用普通 O2 编译，不开启会改变浮点语义的 fast-math；稳定公式也不能恢复读入时已经丢失的间距。

每节写明空集、单点、共线、相切、重合等返回约定；不能把“返回空”统一理解成同一种含义。例如快速半平面交只返回正面积区域，而圆重合有无限交点，需要单独的标记。

\subsubsection*{复制与调用顺序}
每份源码开头列出依赖：先复制整数点或实数点基础，再按依赖顺序加入目标算法，具体文件见各节。二维／三维的 dot、cross、length 等名称通过点类型重载；构造参数或容器须明确类型，避免裸花括号调用的歧义。大多数函数直接传入数据即可，不需要 init 或多测清空；Line、Plane 构造时会单位化方向／法向量，之后须维持此约定。

各节使用教程分别说明接口参数、返回值和退化标记，示例只展示一组数据与调用；代码注释标注函数和局部辅助函数的作用。整数面积与平方量先保留 128 位，确需实数面积／距离时再转换；有限交点、optional、无效圆等必须先检查标记再读取坐标。
